\appendix

\section{Reference Test Vectors}
\label{sec:appendix-vectors}

\paragraph{Purpose.} This appendix lists deterministic test
vectors used by the cross-language (Rust + TypeScript)
reference implementations. Any production or audit
implementation MUST reproduce these outputs bit-equal from the
given inputs. Source: \texttt{crates/tardus-wallet/src/mnemonic.rs}
+ \texttt{ts-sdk/src/mnemonic-rust-compat.ts}.

License: TARDUS-PROPRIETARY-1.0.

\subsection{BIP-39 \texorpdfstring{$\to$}{->} master seed}
\label{sec:vec-master-seed}

The canonical BIP-39 test mnemonic from the original BIP-39
specification (Trezor reference, 12 words):

\begin{center}\small
\begin{tabular}{rl}
\toprule
\texttt{mnemonic}      & \texttt{abandon abandon abandon abandon abandon abandon} \\
                       & \texttt{abandon abandon abandon abandon abandon about}   \\
\texttt{passphrase}    & (empty string)                                            \\
\texttt{master\_seed (TARDUS, first 32 bytes of PBKDF2 seed)} & \\
                       & \texttt{5eb00bbddcf069084889a8ab9155568165f5c453ccb85e70811aaed6f6da5fc1} \\
\bottomrule
\end{tabular}
\end{center}

Reproduce with:
\begin{itemize}[leftmargin=2em,itemsep=0pt]
  \item Rust: \texttt{tardus\_wallet::derive\_master\_seed(\&mnemonic, "")}.
  \item TS: \texttt{deriveMasterSeed(phrase, '')} (v0.1, X25519 path) OR
        \texttt{deriveMasterSeedRustCompat(phrase, '')} (v0.2,
        Rust-compat).
\end{itemize}

\subsection{Master seed \texorpdfstring{$\to$}{->} receiving keypair (Rust + TS v0.2)}
\label{sec:vec-recv-keypair}

\begin{center}\small
\begin{tabular}{rl}
\toprule
\texttt{master\_seed (input)} & \texttt{5eb00bbddcf069084889a8ab9155568165f5c453ccb85e70811aaed6f6da5fc1} \\
\texttt{HKDF salt}            & ASCII \texttt{"TARDUS-recv-id-v1"} \\
\texttt{HKDF info}            & ASCII \texttt{"ed25519-keypair"} \\
\texttt{HKDF output length}   & 64 bytes (HKDF-SHA-512) \\
\texttt{sk (canonical Ed25519 scalar, mod $\ell$)} & \\
                              & \multicolumn{1}{l}{\textit{Deterministic; bit-equal Rust} } \\
                              & \multicolumn{1}{l}{\textit{$\to$ TS v0.2.}} \\
\texttt{pk (32-byte Edwards-compressed point)} & \\
                              & \multicolumn{1}{l}{\textit{$= (sk \cdot G_{\mathrm{ed25519}}).\mathsf{compress}()$.}} \\
\bottomrule
\end{tabular}
\end{center}

Cross-validate by running \texttt{cargo test -p tardus-wallet
test\_canonical\_master\_seed} (Rust) and
\texttt{node --test --experimental-strip-types ts-sdk/test/mnemonic.test.ts}
(TS v0.1) or
\texttt{ts-sdk/test/rust-compat.test.ts} (TS v0.2). The first 20
hex characters of \texttt{master\_seed} are asserted in both
languages' test suites as \texttt{5eb00bbddcf069084889}.

\subsection{Sealed-box AEAD wire format}
\label{sec:vec-sealed-box-wire}

The wire format (32-byte ephemeral public key followed by
AEAD ciphertext + 16-byte Poly1305 tag) is shared between Rust
and TS v0.2.

\begin{center}\small
\begin{tabular}{rl}
\toprule
\texttt{offset 0..32}     & ephemeral public key $epk$ \\
                          & (Curve25519 u-coordinate, little-endian) \\
\texttt{offset 32..(N+48)} & \texttt{ChaCha20-Poly1305}$_{K,N}$\texttt{(plaintext)} \\
                          & where $K$ (32 B) and $N$ (12 B) are HKDF-derived. \\
\bottomrule
\end{tabular}
\end{center}

HKDF inputs (both Rust and TS v0.2 MUST agree byte-for-byte):

\begin{center}\small
\begin{tabular}{rl}
\toprule
\texttt{salt}             & ASCII \texttt{"TARDUS-sealed-box-v1"} (20 bytes) \\
\texttt{ikm (input keying material)} & ECDH shared secret \\
                          & ($= esk \cdot pk_R^{(\mathrm{X25519})}$, 32 bytes) \\
\texttt{info for key}     & ASCII \texttt{"chacha20poly1305-key"} $\,\|\,$ $epk$ $\,\|\,$ $pk_R^{(\mathrm{X25519})}$ \\
\texttt{info for nonce}   & ASCII \texttt{"chacha20poly1305-nonce"} $\,\|\,$ $epk$ $\,\|\,$ $pk_R^{(\mathrm{X25519})}$ \\
\texttt{key output length} & 32 bytes \\
\texttt{nonce output length} & 12 bytes \\
\bottomrule
\end{tabular}
\end{center}

\textbf{TS v0.1 caveat.} The TS v0.1 path (\texttt{src/sealed-box.ts})
uses X25519-clamped scalar derivation per RFC 7748 §5; its
on-the-wire ciphertexts are NOT decryptable by the Rust
reference. The v0.2 path
(\texttt{src/sealed-box-rust-compat.ts}) uses the unclamped
Curve25519 Montgomery ladder
(\texttt{src/montgomery.ts}) and IS wire-byte-compatible.

\subsection{Nullifier computation}
\label{sec:vec-nullifier}

The on-chain nullifier of a coin with public commitment $Cp$:

\[
  \mathsf{nullifier} \;=\; \mathsf{SHA\text{-}256}(\mathtt{"TARDUS-nullifier-v1"} \,\|\, Cp).
\]

Reproduce with
\texttt{tardus\_program::processor::compute\_nullifier(\&cp\_bytes)}.

\subsection{PDA seeds}
\label{sec:vec-pda-seeds}

Solana program-derived addresses for TARDUS PDAs are computed
via \texttt{Pubkey::find\_program\_address(\&[\&str, ...], \&program\_id)}.
Seeds:

\begin{center}\small
\begin{tabular}{rl}
\toprule
\texttt{KeysetRegistry}  & \texttt{["tardus", "keyset-registry"]}                       \\
\texttt{NullifierTree}   & \texttt{["tardus", "nullifier-tree"]}                        \\
\texttt{Vault}           & \texttt{["tardus", "vault", \&denom.to\_le\_bytes()]}          \\
\texttt{SponsorPool}     & \texttt{["tardus", "sponsor-pool"]}                          \\
\bottomrule
\end{tabular}
\end{center}

On devnet with the canonical program ID
\texttt{AmY1ysgQyCC6CmorXkrNkogBSHmomy4kbsPziAYUx47u}, the
KeysetRegistry PDA resolves to
\texttt{4XJ8fmJgJEc9QhNWgNdw1xJHzs3cEZ6JhM746KWDpHK2}.
Reproduce with
\texttt{tardus devnet info --rpc https://api.devnet.solana.com}.

\subsection{Ed25519 precompile envelope}
\label{sec:vec-ed25519-precompile}

The TARDUS \texttt{Refresh} and \texttt{Withdraw} instructions
require an immediately-preceding ed25519 precompile instruction
with the following 144-byte data layout (single-signature
variant):

\begin{center}\small
\begin{tabular}{rl}
\toprule
\texttt{offset 0..1}    & \texttt{num\_signatures = 1}      \\
\texttt{offset 1..2}    & \texttt{padding = 0}              \\
\texttt{offset 2..4}    & \texttt{sig\_offset = 16} (u16 LE) \\
\texttt{offset 4..6}    & \texttt{sig\_ix\_idx = 0xFFFF} (u16 LE) \\
\texttt{offset 6..8}    & \texttt{pk\_offset = 80}          \\
\texttt{offset 8..10}   & \texttt{pk\_ix\_idx = 0xFFFF}     \\
\texttt{offset 10..12}  & \texttt{msg\_offset = 112}        \\
\texttt{offset 12..14}  & \texttt{msg\_size = 32}           \\
\texttt{offset 14..16}  & \texttt{msg\_ix\_idx = 0xFFFF}    \\
\texttt{offset 16..80}  & 64-byte signature $\sigma$ on coin $Cp$ \\
\texttt{offset 80..112} & 32-byte mint joint pubkey $\mathsf{joint\_pk}$ \\
\texttt{offset 112..144} & 32-byte coin commitment $Cp$ (the message) \\
\bottomrule
\end{tabular}
\end{center}

Reproduce with
\texttt{crates/tardus-cli/src/commands/devnet.rs::build\_ed25519\_precompile\_data}
or
\texttt{crates/tardus-wallet-gui/src/runtime.rs::withdraw\_on\_devnet}.

\subsection{\texorpdfstring{Solana realloc 10\,KB per-instruction cap (discovery \#10)}{Solana realloc 10 KB per-instruction cap (discovery 10)}}
\label{sec:vec-realloc-cap}

\texttt{AccountInfo::realloc(new\_size, true)} is bounded by
\texttt{MAX\_PERMITTED\_DATA\_INCREASE = 10240 bytes} per single
instruction call. Growing a PDA from 1024 bytes to 65536 bytes
therefore requires
$\lceil (65536 - 1024) / 10240 \rceil = 7$ sequential
\texttt{ResizeAccount} instructions, each within its own
transaction (or chained in a single transaction with multiple
program calls). The CLI's
\texttt{tardus devnet resize-nullifier-tree} subcommand handles
the multi-transaction batching transparently; operators see a
single high-level target size.
